\section{Exercice 2.1}
\subsection{Consignes de l'exercice}
\begin{tcb-instructions}[title=Exercice 2.1]
A l’aide de \LaTeX{ }et de votre éditeur de texte préféré :

    \begin{itemize}[label=\checkmark]
        \item \textbf{Réalisez un document au format PDF de 4 pages au moins} comportant : une page de titre (avec vos prénom(s) et nom(s)), une Table des Matières, un titre (ou plusieurs) de section (sous-section...), plusieurs paragraphes (vous pouvez reprendre les éléments informationnels de votre dossier de presse, ou choisir une autre thématique), avec \textbf{au moins} : une image, un bout de code avec coloration syntaxique du langage de votre choix, une référence bibliographique (avec ou sans \textbf{BibTeX} selon) et un tableau.
        \item Lorsque votre PDF est édité, signez-le numériquement de la manière de votre choix.
        \item Ajoutez ensuite un \textbf{graphique TIKZ} à votre document LATEX.
        \item Les fans de Raoul Capablanca, de Magnus Carlsen ou du \textit{Jeu de la Dame} pourront dessiner un ou plusieurs échiquiers pour expliquer une attaque / défense de leur choix à la place du graphique TikZ.
    \end{itemize}

\textbf{Important :} Les étudiant.e.s qui se sont inscrits à l’un des cours suivants : ADO, IF-UOR, RES, SE ou I\&C (ancienne maquette) ou HTML \& CSS, Linux, OMD (ce cours), I\&C (nouvelle maquette) et qui ont déjà intégré \LaTeX{ }pour leurs rendus d’exercices et PDF n’ont pas besoin de réaliser cet exercices 2.1 (sauf s’ils/si elles le souhaitent). Il en sera de même pour tous les étudiant.e.s qui réaliseront le rendu des TP(s) de ce cours en PDF généré à partir de \LaTeX{ }. Il vous suffit de le préciser dans la réponse et d’ajouter un C.C de votre source \LaTeX{ }(du chapitre de leur choix) et d’insérer le fichier \path{*.tex} dans l’archive attachée (+ dépôt dans votre espace \textit{GitHub}).

\end{tcb-instructions}

\subsection{Réalisation de l'exercice}
Ayant réalisé l'intégralité de ce document grâce à \textit{TexLive} et \textit{TexStudio}, je place donc le copier-coller de l'intégralité du code de ce chapitre en annexe (\autoref{ann:ex-2.1_code-latex}, \autopageref{ann:ex-2.1_code-latex}). J'évite les autres chapitres car ils sont bien plus longs et il me semble inutile de copier-coller autant de lignes...

Par acquis de conscience, en plus du fichier \path{chapitre-2.tex}, je place aussi le fichier \path{chapitre-1.tex} ainsi que mes fichiers de configuration (\path{main.tex}, \path{packages.tex}, \path{environments.tex}, \path{colours.tex}, \path{titlepage.tex} et \path{metadata.tex}) dans l'archive, dans l'état dans lequel ils sont au moment de la rédaction de cette réponse (au 21 novembre 2026 donc).

\subsection{Conclusion}
Je n'avais jamais eu l'occasion d'essayer \LaTeX{ }dans ma vie professionnelle malgré tout le bien que l'on m'en avait dit (j'ai toujours utilisé \textit{LibreOffice}) et je suis très agréablement surpris par la quantité de temps que libère la gestion automatique de la mise en page.

Il m'a bien fallu deux jours à temps complet pour comprendre le principe et créer un modèle qui soit à la fois simple et dont je maitrise tout le contenu, mais c'est un investissement qui me fait gagner beaucoup de temps aujourd'hui!